\chapter{Heat Exhaustion}

\section*{Definition}
Heat exhaustion is a moderate heat-related illness characterized by heavy sweating, weakness, dizziness, and nausea due to inadequate fluid and electrolyte replacement during heat exposure. It can progress to heat stroke if not treated promptly.

\section*{What Happens in Your Body}
Prolonged exposure to high environmental temperatures causes excessive sweating with loss of water and electrolytes. The body attempts to maintain thermoregulation through widening of blood vessels and increased cardiac output, but volume depletion leads to reduced stroke volume and cardiac output, causing fatigue, weakness, dizziness, and orthostatic hypotension.

\section*{Causes}
\begin{itemize}
\item Strenuous exercise in hot, humid conditions
\item Prolonged exposure to high temperatures
\item Inadequate fluid intake
\item Excessive sweating with inadequate salt replacement
\item Heat waves affecting vulnerable populations
\item Medications that impair thermoregulation (diuretics, beta-blockers)
\item Alcohol consumption
\item Obesity
\item Advancing age
\end{itemize}

\section*{When to See a Doctor}
\begin{itemize}
\item Heavy sweating, cool and clammy skin
\item Weakness and fatigue
\item Dizziness or lightheadedness
\item Nausea and possible vomiting
\item Headache
\item Mild confusion or irritability
\item Thirst
\item Decreased urine output
\item Body temperature normal to mildly elevated (37-40 C)
\item Orthostatic hypotension, rapid pulse
\end{itemize}

\section*{Related Symptoms}
\begin{itemize}
\item Heat stroke (severe progression)
\item Heat cramps
\item Heat syncope (fainting)
\item Dehydration
\item Hyponatremia (low sodium)
\item Hypovolemia
\end{itemize}

\section*{Self-Care / Home Management}
\begin{itemize}
\item Move to a cool, shaded or air-conditioned environment
\item Lie down and elevate the feet
\item Drink cool water or electrolyte-containing sports drinks
\item Remove tight or excess clothing
\item Apply cool, wet cloths or take a cool shower
\item Rest for 24 hours after recovery
\item Avoid strenuous activity and heat exposure for several days
\end{itemize}

\section*{How Your Doctor Will Evaluate You}
\begin{itemize}
\item Clinical diagnosis based on history of heat exposure and symptoms
\item Core body temperature measurement
\item Electrolytes, renal function, BUN-to-creatinine ratio
\item Urinalysis (specific gravity, color)
\end{itemize}

\section*{Treatment Options}
\begin{itemize}
\item Rest in a cool environment
\item Oral or IV fluid replacement with balanced electrolyte solutions
\item Monitor for progression to heat stroke
\item Avoid exertion and heat until fully recovered
\item Education on prevention and hydration strategies
\end{itemize}
\section*{References}
\begin{thebibliography}{99}
\bibitem{heat_exhaustion:ref1} Bouchama A, Knochel JP. "Heat Stroke." N Engl J Med. 2002;346(25):1978--1988.
\bibitem{heat_exhaustion:ref2} Casa DJ, DeMartini JK, Bergeron MF, et al. "National Athletic Trainers' Association Position Statement: Exertional Heat Illnesses." J Athl Train. 2015;50(9):986--1000.
\bibitem{heat_exhaustion:ref3} Lipman GS, Eifling KP, Ellis MA, et al. "Wilderness Medical Society Practice Guidelines for the Prevention and Treatment of Heat-Related Illness." Wilderness Environ Med. 2013;24(4):351--361.
\bibitem{heat_exhaustion:ref4} Yeo TP. "Heat Stroke: A Comprehensive Review." AACN Clin Issues. 2004;15(2):280--293.
\bibitem{heat_exhaustion:ref5} Glazer JL. "Management of Heatstroke and Heat Exhaustion." Am Fam Physician. 2005;71(11):2133--2140.
\end{thebibliography}
